\section*{Capítulo \ref{ch:b3:chromatin-epigenetics} --- Cromatina e epigenética}

\begin{solution}{exo:b3:chromatin-epigenetics:1}
Duas cópias de cada uma das \omterm{def:b3:chromatin-epigenetics:nucleosome}{histonas} H2A, H2B, H3 e H4 (um octâmero)
com \qty{147}{bp} de DNA em $1.65$ voltas; ligantes de
\qtyrange{20}{60}{bp}, o que dá uma repetição próxima de \qty{200}{bp}.
A nuclease corta apenas o DNA ligante exposto, nunca dentro de um
cerne, de modo que todo fragmento é um número inteiro de repetições:
uma escada de \num{200}, \num{400}, \qty{600}{bp}, e não um rastro.
\end{solution}

\begin{solution}{exo:b3:chromatin-epigenetics:2}
Diploide: $6.0\times 10^{9}$ bp; $6.0\times 10^{9}/190 \approx
3.2\times 10^{7}$ \omterm{def:b3:chromatin-epigenetics:nucleosome}{nucleossomos}; comprimento $6.0\times 10^{9}\times
\qty{0.34}{nm} = \qty{2.0}{m}$.
\end{solution}

\begin{solution}{exo:b3:chromatin-epigenetics:3}
H3K27me3: \omterm{def:b3:chromatin-epigenetics:nucleosome}{histona} H3, lisina 27, trimetilada --- \omterm{def:b3:chromatin-epigenetics:nucleosome}{cromatina} silenciosa
de genes reprimidos no desenvolvimento, escrita pelo PRC2 e lida pelas
proteínas com cromodomínio do PRC1. H4K16ac: H4, lisina 16, acetilada
--- \omterm{def:b3:chromatin-epigenetics:nucleosome}{cromatina} ativa e aberta (rompe um contato
nucleossomo--nucleossomo). H3S10ph: H3, serina 10, fosforilada --- uma
marca mitótica associada à condensação cromossômica (e,
transitoriamente, à ativação gênica por alguns sinais); não é uma
marca estável de silenciamento nem de ativação.
\end{solution}

\begin{solution}{exo:b3:chromatin-epigenetics:4}
O gene alaranjado é ligado ao X e tem um alelo alaranjado e um preto;
uma pelagem malhada exige um de cada, em dois cromossomos X
diferentes, um dos quais é silenciado por célula. Um macho tem um só X
e é inteiramente alaranjado ou inteiramente preto. Um macho tricolor
carrega dois cromossomos X e um Y (XXY), e é estéril.
\end{solution}

\begin{solution}{exo:b3:chromatin-epigenetics:5}
$\qty{200}{bp}\times \qty{0.34}{nm} = \qty{68}{nm}$ por conta de
\qty{11}{nm}: razão $6.2$. Necessária: $\qty{4}{cm}/\qty{4}{\micro m} =
10^{4}$. Dobramento adicional de $10^{4}/6.2 \approx 1600$.
\end{solution}

\begin{solution}{exo:b3:chromatin-epigenetics:6}
$\mu = 0.98$, $\delta = 0.01$: $p^{*} = 0.01/0.03 = 0.33$; razão $\mu -
\delta = 0.97$; meia-vida $\ln 2/(-\ln 0.97) = 22.8 \approx 23$
divisões. $\mu = 0.999$, $\delta = 0.001$: $p^{*} = 0.5$; razão
$0.998$; meia-vida $\ln 2 / 0.002 \approx 350$ divisões.
\end{solution}

\begin{solution}{exo:b3:chromatin-epigenetics:7}
$A$: ativo (promotor com H3K4me3 e \omterm{def:b3:chromatin-epigenetics:histone-code}{acetilação}). $B$: silencioso,
reprimido pelo Polycomb --- \omterm{def:b3:chromatin-epigenetics:nucleosome}{heterocromatina} facultativa, o estado de um
gene do desenvolvimento desligado nesta linhagem mas usado em outras,
de modo que $B$ é o que tem mais chance de estar ativo em outro tipo
celular. $C$: silencioso, \omterm{def:b3:chromatin-epigenetics:nucleosome}{heterocromatina} constitutiva (via da HP1),
tipicamente repetições ou genes silenciados em todos os tipos
celulares.
\end{solution}

\begin{solution}{exo:b3:chromatin-epigenetics:8}
O \emph{UBE3A} é expresso apenas a partir do alelo materno nos
neurônios. A deleção dela veio do pai, no alelo que já está silencioso
de qualquer modo, e por isso ela é saudável; quando ela a transmite, a
deleção passa a ser materna: cada filho tem probabilidade $1/2$ de
herdá-la e então tem a síndrome de Angelman. Se a deleção tivesse vindo
da mãe dela, ela própria teria a síndrome de Angelman; uma portadora
saudável só pode tê-la herdado por via paterna. (Uma deleção paterna de
toda a região daria Prader--Willi; o \emph{UBE3A} sozinho não basta
para isso.)
\end{solution}

\begin{solution}{exo:b3:chromatin-epigenetics:9}
Um X sem \emph{\omterm{def:b3:chromatin-epigenetics:x-inactivation}{Xist}} não pode ser inativado, de modo que o outro X é
silenciado em todas as células: a inativação é completa, mas deixa de
ser aleatória (é enviesada), e o embrião é viável. Os dois X sem
\emph{\omterm{def:b3:chromatin-epigenetics:x-inactivation}{Xist}}: nenhuma inativação, dois cromossomos X ativos, expressão
dobrada dos genes ligados ao X, e o embrião fêmea morre cedo. Um
transgene \emph{\omterm{def:b3:chromatin-epigenetics:x-inactivation}{Xist}} num autossomo recobre esse autossomo em cis e
silencia seus genes, causando perda de função de uma cópia autossômica
--- o experimento que mostrou que o \emph{\omterm{def:b3:chromatin-epigenetics:x-inactivation}{Xist}} basta para silenciar
em cis.
\end{solution}

\begin{solution}{exo:b3:chromatin-epigenetics:10}
As citosinas metiladas se desaminam a timina, que o reparo não
reconhece, de modo que um \omterm{def:b3:chromatin-epigenetics:methylation}{CpG} metilado muta para TpG ou CpA a uma taxa
alta. Só as sequências não metiladas na \emph{linhagem germinativa},
geração após geração, conservam seus \omterm{def:b3:chromatin-epigenetics:methylation}{CpGs}; as ilhas conservaram os
seus, e portanto estiveram não metiladas na linhagem germinativa
durante a maior parte da história dos vertebrados. Um gene metilado
apenas no fígado está não metilado no espermatozoide e no óvulo; as
mutações que surgem nas células do fígado não são herdadas, de modo que
a metilação somática não deixa pegada na sequência.
\end{solution}

\begin{solution}{exo:b3:chromatin-epigenetics:11}
Na replicação as \omterm{def:b3:chromatin-epigenetics:nucleosome}{histonas} antigas, marcadas, são repartidas entre os
dois dúplexes-filhos e intercaladas com \omterm{def:b3:chromatin-epigenetics:nucleosome}{histonas} novas sem marca, de
modo que cada domínio-filho fica cerca de metade marcado. A HP1 ligada
a um \omterm{def:b3:chromatin-epigenetics:nucleosome}{nucleossomo} marcado retido recruta a SUV39H, que metila os
vizinhos sem marca; enquanto os \omterm{def:b3:chromatin-epigenetics:nucleosome}{nucleossomos} marcados forem densos o
bastante para levar um \omterm{def:b3:chromatin-epigenetics:histone-code}{escritor} a cada lacuna, o domínio é recomposto
antes da divisão seguinte --- a mesma retroalimentação do $\mu$ alto do
teorema. A propagação para nas fronteiras: proteínas isolantes (CTCF),
genes muito transcritos e genes de tRNA, \omterm{def:b3:chromatin-epigenetics:nucleosome}{nucleossomos} com marcas ativas
incompatíveis (H3K4me3, \omterm{def:b3:chromatin-epigenetics:histone-code}{acetilação}) que o \omterm{def:b3:chromatin-epigenetics:histone-code}{escritor} não consegue usar, e
o ponto em que os \omterm{def:b3:chromatin-epigenetics:histone-code}{apagadores} (desmetilases, renovação das HDAC)
superam o \omterm{def:b3:chromatin-epigenetics:histone-code}{escritor}. Teste: recrute a HP1 transitoriamente para um gene
repórter (uma fusão com um domínio de ligação ao DNA cuja ligação seja
controlada por fármaco), solte-a e acompanhe o silêncio do repórter e a
H3K9me3 ao longo das divisões; preveja uma persistência que depende do
\omterm{def:b3:chromatin-epigenetics:histone-code}{escritor}, e a perda se o \omterm{def:b3:chromatin-epigenetics:histone-code}{escritor} ou a dimerização do \omterm{def:b3:chromatin-epigenetics:histone-code}{leitor} estiverem
mutados, ou se um isolante for colocado dentro do domínio.
\end{solution}

\begin{solution}{exo:b3:chromatin-epigenetics:12}
Excluir: um ambiente compartilhado (a dieta e a pobreza dos próprios
netos); um efeito uterino --- a filha da mãe exposta à fome era um
feto, com suas células germinativas já presentes, de modo que a
linhagem germinativa do neto foi exposta diretamente e o efeito não é
transgeracional; diferenças genéticas entre as famílias que
sobreviveram e as que não; causalidade reversa (a própria massa
corporal altera a metilação no sangue); diferenças na composição das
células do sangue; e testes múltiplos ao longo de muitos \omterm{def:b3:chromatin-epigenetics:methylation}{CpGs}.
Experimento em camundongos: exponha \emph{machos} F$_0$ de linhagem
isogênica (sem via uterina), cruze-os com fêmeas não expostas ou use
fecundação in vitro e transferência de embriões, e acompanhe F$_2$ e
F$_3$ pela linha paterna, medindo a metilação no esperma de cada
geração; uma marca que persiste no esperma e no fenótipo da F$_3$, numa
linhagem sem variação genética, é transmissão germinativa.
\end{solution}

\begin{solution}{pb:b3:chromatin-epigenetics:1}
\textbf{1.} $6.4\times 10^{9}\times \qty{0.34}{nm} = \qty{2.2}{m}$;
$6.4\times 10^{9}/200 = 3.2\times 10^{7}$ \omterm{def:b3:chromatin-epigenetics:nucleosome}{nucleossomos}.
\textbf{2.} \omterm{def:b3:chromatin-epigenetics:nucleosome}{Histonas}: $3.2\times 10^{7}\times 1.08\times 10^{5} =
3.5\times 10^{12}$ Da $= \qty{5.7}{pg}$. DNA: $6.4\times 10^{9}\times
650 = 4.2\times 10^{12}$ Da $= \qty{6.9}{pg}$. Massas comparáveis.
\textbf{3.} Núcleo: $\tfrac{4}{3}\pi\,(\qty{5}{\micro m})^{3} =
\qty{520}{\micro m^3}$. Cerne: $\pi\,(5.5)^{2}\times 5.5 =
\qty{520}{nm^3} = 5.2\times 10^{-7}\,\unit{\micro m^3}$; total
$3.2\times 10^{7}\times 5.2\times 10^{-7} = \qty{17}{\micro m^3}$:
cerca de \qty{3}{\%} do núcleo.
\textbf{4.} $3.2\times 10^{7}\times \qty{11}{nm} = \qty{0.35}{m}$;
\omterm{prop:b3:chromatin-epigenetics:compaction}{razão de empacotamento} $2.2/0.35 \approx 6$.
\textbf{5.} $6.4\times 10^{9}/2\times 10^{5} = \num{32000}$ alças de
\num{1000} \omterm{def:b3:chromatin-epigenetics:nucleosome}{nucleossomos} cada.
\textbf{6.} $0.21^{2}\times 6.4\times 10^{9} = 2.8\times 10^{8}$;
observados $5.6\times 10^{7}$, ou seja \qty{80}{\%} perdidos, por
desaminação da 5-metilcitosina a timina (transições C$\to$T não
reparadas nos \omterm{def:b3:chromatin-epigenetics:methylation}{CpGs} metilados).
\textbf{7.} $p_{n+1} = 0.99\,p_{n} + 0.02\,(1 - p_{n}) = 0.97\,p_{n} +
0.02$; $p^{*} = 0.02/0.03 = 0.67$.
\textbf{8.} $p_{n} = 0.67 + 0.33\times 0.97^{n}$: $n = 10$: $0.91$;
$n = 50$: $0.74$; $n = 200$: $0.67$ (o excesso é $7\times 10^{-4}$).
\textbf{9.} $\ln 2/(-\ln 0.97) = 22.8 \approx 23$ divisões; a uma por
semana, cerca de \num{23} semanas --- cinco meses.
\textbf{10.} Cada replicação emparelha uma fita antiga, ainda metilada,
com uma fita nova que enzima nenhuma metila; as fitas antigas não são
nem ganhas nem perdidas, o número de fitas dobra, e portanto a fração
metilada cai pela metade. $0.80/2^{n} < 0.05$ exige $2^{n} > 16$: cinco
divisões ($0.025$; quatro dão exatamente \qty{5}{\%}).
\textbf{11.} Razão $0.9995 - 0.0005 = 0.999$; meia-vida $\ln
2/(-\ln 0.999) \approx 690$ divisões, cerca de \num{13} anos a uma por
semana.
\textbf{12.} Um gene silencioso por uma vida inteira de umas \num{4000}
divisões semanais precisa de uma meia-vida de centenas de divisões, que
só a retroalimentação da questão~11 dá: os \omterm{def:b3:chromatin-epigenetics:histone-code}{leitores} da marca (MeCP2,
HP1) recrutam seus \omterm{def:b3:chromatin-epigenetics:histone-code}{escritores} (metilase da H3K9, DNMT3), de modo que um
domínio denso de marcas se copia com uma eficiência que enzima nenhuma
tem sozinha.
\textbf{13.} Um X carrega o alelo alaranjado, o outro o preto; toda
célula da pele silenciou um X e expressa apenas o outro; os
descendentes de uma célula mantêm sua escolha, de modo que uma mancha é
um clone (ou um grupo de clones vizinhos) que fez a mesma escolha.
\textbf{14.} Todas as $N$ escolhem o X que carrega o alaranjado, ou
todas escolhem o outro: $2\times (1/2)^{N} = 2^{1-N}$.
\textbf{15.} $2^{1-N} = 10^{-9}$: $N - 1 = \log_{2} 10^{9} = 29.9$,
$N \approx 30$ células.
\textbf{16.} $\qty{3000}{cm^2}/\qty{15}{cm^2} = 200$ manchas, bem mais
que $30$ fundadoras: durante o crescimento os descendentes de uma
fundadora são dispersados pela mistura e pela migração celular, de modo
que um clone forma várias ilhas (e vizinhas com a mesma escolha se
fundem, o que puxa no sentido contrário, mas menos).
\textbf{17.} Um \omterm{def:b3:chromatin-epigenetics:x-inactivation}{corpúsculo de Barr} por X além do primeiro: XXY um, XXX
dois, XY nenhum.
\textbf{18.} Tolerado quando o gene tem um homólogo funcional no Y, de
modo que os machos também têm duas cópias (os genes pseudoautossômicos
e pares como \emph{KDM5C}/\emph{KDM5D}); importa nas aneuploidias do X
--- na síndrome de Turner (XO) os genes que escapam estão presentes em
cópia única, que é por que um indivíduo XO não é simplesmente uma fêmea
normal (haploinsuficiência do \emph{SHOX} e baixa estatura), e no XXY
os que escapam ficam em dose excessiva.
\textbf{19.} $1/\num{15000} = 6.7\times 10^{-5}$; deleções $0.70\times
6.7\times 10^{-5} = 4.7\times 10^{-5}$ (uma em \num{21000}); dissomia
materna $1.7\times 10^{-5}$ (uma em \num{60000}).
\textbf{20.} A síndrome de Angelman é a perda de um só gene, o
\emph{UBE3A}, de modo que uma única mutação pontual no alelo materno
basta. Prader--Willi é a perda de vários genes de expressão paterna ao
mesmo tempo (\emph{SNRPN}, o agrupamento de RNAs pequenos); uma mutação
pontual inativa um deles e dá no máximo um fenótipo parcial, nunca a
síndrome inteira.
\textbf{21.} \qty{20}{\%} dos filhos recebem uma deleção paterna e têm
Prader--Willi; nenhum tem Angelman, que exige uma deleção materna.
\textbf{22.} Três cromossomos, dois maternos e um paterno, um perdido
ao acaso: o paterno se perde com probabilidade $1/3$, o que deixa
dissomia materna e Prader--Willi; com probabilidade $2/3$ perde-se uma
cópia materna e a criança é normal.
\textbf{23.} Os genes paternos deveriam promover a demanda sobre a mãe;
sua perda (Prader--Willi) deveria reduzir a demanda --- o que se encaixa
na sucção fraca e no baixo ganho de peso do recém-nascido ---, ao passo
que o apetite insaciável posterior é o oposto e é mais difícil de
encaixar (uma leitura proposta é que os genes paternos normalmente
empurram a criança rumo ao desmame). Angelman, a perda de um gene
materno, deveria aumentar a demanda; a sucção prolongada e os despertares
noturnos frequentes das crianças com Angelman se encaixam, as convulsões
e a perda da fala não têm nada a ver com recursos.
\textbf{24.} Ter como alvo o transcrito antissenso paterno
(\emph{UBE3A-ATS}) com um oligonucleotídeo antissenso que o degrade ou
bloqueie sua elongação: o \emph{UBE3A} paterno está intacto e é
silenciado apenas em cis por esse RNA, de modo que retirar o RNA
restabelece sua expressão. A metilação da \omterm{def:b3:chromatin-epigenetics:imprinting}{região controladora do
imprinting} fica intocada, e portanto o \emph{SNRPN} e os demais genes
sob imprinting conservam seu padrão parental normal.
\textbf{25.} Meia-vida de uma marca: cerca de $23$ divisões sem
retroalimentação, cerca de $690$ com ela; a pelagem tricolor foi
decidida por uma população fundadora de cerca de $30$ células.
\end{solution}
